\itemhead{习题1.1 · 第1题（教材p.12）}{XJ-B1-C01-HW-P012-Y0440}
\label{XJ-B1-C01-HW-P012-Y0440}
{\footnotesize 湘教版教材p.12 / PDF第19页 / 习题1.1 / 第1题}
\field{原题}{判断下列各组对象能否构成集合。若能构成集合，指出是有限集还是无限集；若不能构成集合，试说明理由。
\subq{1}我国古代四大名著的名称；
\subq{2}尾数是 5 的自然数；
\subq{3}我们班身高大于 $1.7\,\mathrm m$ 的同学。}
\field{参考答案}{均能构成集合；（1）（3）为有限集，（2）为无限集。判断标准都是确定的。}
{\small \field{关联知识}{S01 集合、元素与常用数集（教材帮p.1、2 / PDF10、11）}
\field{方法与易错点}{先辨认研究对象和数集范围，再代入定义判断；含参集合逐个来源分类。 易错：把自然数集误作正整数集；混淆元素和单元素集合；候选参数造成元素重复。}
\field{答案依据}{答案册相关题组：PDF第4、5页（册内页3、4）；按当前题库原题逐项重算，不按旧版题号机械对套。}}
